%------------------------------------------------------------------------------%
\begin{question}
  Escreva os números complexos na sua forma polar
  \begin{enumerate}
    \item \(z_1 = \sqrt{3} + i\)
    \item \(z_2 = -2 + 2\sqrt{3}i\)
  \end{enumerate}
\end{question}

%------------------------------------------------------------------------------%
\begin{resolution}{Solução}
  \[
    z_1 = \sqrt{3} + i
    \qquad\qquad
    \pause
    a = \Re(z) = \sqrt{3}
    \qquad\qquad
    \pause
    b = \Im(z) = 1
  \]

  \[
    \pause
    \rho
    = \sqrt{a^2 + b^2}
    \pause
    = \sqrt{\left(\sqrt{3}\right)^2 + 1^2}
    \pause
    = \sqrt{3+1}
    \pause
    = \sqrt{4}
    \pause
    = 2
  \]

  \[
    \pause
    \sin(\varphi)
    \pause
    = \frac{b}{\rho}
    \pause
    = \frac{1}{2}
    \pause
    \qquad\qquad
    \cos(\varphi)
    \pause
    = \frac{a}{\rho}
    \pause
    = \frac{\sqrt{3}}{2}
  \]

  \[
    \pause
    \varphi
    \pause
    = \ang{30}
    \pause
    = \frac{\pi}{6} \si{\radian}
  \]

  \[
    \pause
    z_1
    \pause
    = \rho\left(\cos(\varphi)+i\sin(\varphi)\right)
    \pause
    = 2\left(\cos\left(\frac{\pi}{6}\right)+i\sin\left(\frac{\pi}{6}\right)\right)
  \]
\end{resolution}

%------------------------------------------------------------------------------%
